% !PRINTED 15
\weberpage{15}

%% [running head, not body text -- commented out 2026-08-07] Page furniture (running
%% head, folio, signature) is excluded from this edition by standing policy.
%% \gesperrt{\S.\,5.} \hfill \gesperrt{Mehrstufiger Isomorphismus.} \hfill 15

%<<CONT
denote it, following the example of Hölder\footnote{Mathematische Annalen, Bd.\,34. The symbol recalls a quotient, with which the complementary group does indeed have a certain analogy.}, by $P/N$. The degree of the complementary group is equal to the index of the group $N$.

At the close of these considerations we mention one further theorem on normal divisors.

\noindent 5.\ \gesperrt{If $N$ is a divisor of $P$, and $N'$ a divisor of $N$ and at the same time a normal divisor of $P$, then $N'$ is also a normal divisor of $N$.}

This is self-evident, because for every element $b$ of $P$ we have $b^{-1}N'b = N'$, and because every element of $N$ is also an element of $P$.

\gesperrt{But one may not conclude conversely that, if $N$ is a normal divisor of $P$ and $N'$ a normal divisor of $N$, then $N'$ must also be a normal divisor of $P$.}

\begin{center}
\S.\,5.

\gesperrt{Multiple Isomorphism.}
\end{center}

If, as above, $N$ is a normal divisor of $P$ of degree $\nu$ and index $\mu$, so that $n = \mu \nu$, then the group $P/N$ complementary to $P$ is of degree $\mu$. We will denote this group, or one isomorphic with it, by $Q$, and its elements, which correspond to the subsidiary groups $N, Nb_1, Nb_2 \dots$, by $A, B, C \dots$ To each of these elements there correspond $\nu$ elements of the group $P$, say

\noindent to the element $A$ the elements $a, a_1 \dots a_{\nu-1}$

\noindent to the element $B$ the elements $b, b_1 \dots b_{\nu-1}$

\noindent $. \; . \; . \; . \; . \; . \; . \; . \; . \; . \; . \; . \; . \; . \; . \; . \; . \; . \; . \; . \; .$

This correspondence is then of such a kind that every composed element $ab$ corresponds to the composed element $AB$. For the elements $A$ and $B$ correspond to the subsidiary groups $Na$ and $Nb$, and, because $N$ is a normal divisor, we have
\[
Na \, Nb = Nab.
\]

Thus in the composition of the $A, B \dots$ on the one hand and of the $a, b \dots$ on the other, the same law holds here as in
%CONT>>
